
\documentclass{memoir}

\usepackage{sourcecodepro}
\usepackage{sourcesanspro}
\usepackage{sourceserifpro}

\usepackage{fancyvrb}
\usepackage{fvextra}
\usepackage{xparse}

\usepackage[most]{tcolorbox}

% Detect whether we're in a footnote. This is used later to avoid \href rendering bugs.
\newif\ifinfootnote
\infootnotefalse

\let\oldfootnote\footnote
\renewcommand{\footnote}[1]{%
  \oldfootnote{\infootnotetrue#1\infootnotefalse}%
}

\let\oldfootnotetext\footnotetext
\renewcommand{\footnotetext}[1]{%
  \oldfootnotetext{\infootnotetrue#1\infootnotefalse}%
}

\usepackage[breaklinks=true, hyperfootnotes=false]{hyperref}

% Redefines \href to show footnotes with URLs instead. This works around a page breaking bug in
% hyperref and also makes the link useful on paper. If already in a footnote, the URL is
% in parentheses instead.
\let\oldhref\href
\RenewDocumentCommand{\href}{mm}{%
  \ifinfootnote%
    #2~(\url{#1})%
  \else%
    #2\footnote{\url{#1}}%
  \fi%
}

\usepackage[normalem]{ulem}
\newcommand{\coloredwave}[2]{\textcolor{#1}{\uwave{\textcolor{black}{#2}}}}
\usepackage{newunicodechar}

% Work around the fact that
% U+271D LATIN CROSS doesn't exist in
% DejaVu Sans Mono Oblique. \textup
% is fontspec for "upright, not italic/oblique".
\newunicodechar{✝}{\textup{✝}}
% Work around missing U+2011 (non-breaking hyphen) in Source Serif Pro
\newunicodechar{‑}{-}

\definecolor{errorColor}{HTML}{B91C1C}
\definecolor{infoColor}{HTML}{1E6BB8}
\definecolor{warningColor}{HTML}{D97706}
\newcommand{\errorDecorate}[1]{\coloredwave{errorColor}{#1}}
\newcommand{\infoDecorate}[1]{\coloredwave{infoColor}{#1}}
\newcommand{\warningDecorate}[1]{\coloredwave{warningColor}{#1}}
\DefineVerbatimEnvironment{LeanVerbatim}{Verbatim}
  {commandchars=\\\{\},fontsize=\small,breaklines=true}
\DefineVerbatimEnvironment{FileVerbatim}{Verbatim}{commandchars=\\\{\},fontsize=\small,breaklines=true,frame=single,framesep=2mm,numbers=left}
\CustomVerbatimCommand{\LeanVerb}{Verb}
  {commandchars=\\\{\},fontsize=\small,breaklines=true}
\CustomVerbatimCommand{\FileListingVerb}{Verb}
  {commandchars=\\\{\},fontsize=\small,frame=single,framesep=2mm, numbers=left}

%%% Trace messages
\newlength{\traceindent}
\setlength{\traceindent}{1.5em}

\newcommand{\expandedIndicator}{$\blacktriangledown$\hspace{2pt}}
\newcommand{\collapsedIndicator}{$\blacktriangleright$\hspace{2pt}}

\newenvironment{expandedtrace}[1]{%
  \par\noindent\expandedIndicator #1\par
  \advance\leftskip by \traceindent
}{%
}

\newenvironment{collapsedtrace}[1]{%
  \par\noindent\collapsedIndicator #1\par
  \advance\leftskip by \traceindent
}{%
}
%%% End Trace messages

\definecolor{bordercolor}{HTML}{98B2C0}
\definecolor{medgray}{HTML}{555555}
\newtcolorbox{docstringBox}[2][]{colback=white,
breakable,
colframe=bordercolor,
colbacktitle=white,
enhanced,
coltitle=medgray,
attach boxed title to top left={xshift=2mm,yshift=-2mm},
boxrule=0.4pt,
fonttitle=\sffamily\fontsize{6pt}{7pt}\selectfont,
boxed title style={top=-0.3mm,bottom=-0.3mm,left=-0.3mm,right=-0.3mm,boxrule=0.4pt},
title={#2},#1}


\makechapterstyle{lean}{%
\renewcommand*{\chaptitlefont}{\sffamily\HUGE}
\renewcommand*{\chapnumfont}{\chaptitlefont}
% allow for 99 chapters!
\settowidth{\chapindent}{\chapnumfont 999}
\renewcommand*{\printchaptername}{}
\renewcommand*{\chapternamenum}{}
\renewcommand*{\chapnumfont}{\chaptitlefont}
\renewcommand*{\printchapternum}{%
\noindent\llap{\makebox[\chapindent][l]{%
\chapnumfont \thechapter}}}
\renewcommand*{\afterchapternum}{}
}

\chapterstyle{lean}

\setsecheadstyle{\sffamily\bfseries\Large}
\setsubsecheadstyle{\sffamily\bfseries\large}
\setsubsubsecheadstyle{\sffamily\bfseries}

\renewcommand{\cftchapterfont}{\normalfont\sffamily}
\renewcommand{\cftsectionfont}{\normalfont\sffamily}
\renewcommand{\cftchapterpagefont}{\normalfont\sffamily}
\renewcommand{\cftsectionpagefont}{\normalfont\sffamily}
\setmonofont{DejaVu Sans Mono}

\title{\sffamily Capability Monotonicity in a Formal Skeleton for Agentic AI Safety}
\author{\sffamily Manish Bhatt}
\date{\sffamily }

% --- added by build-ctvg-thesis.sh: glyph coverage and line breaking
\newunicodechar{∎}{\leavevmode\rule{0.5em}{0.5em}}
\newunicodechar{ε}{\ensuremath{\varepsilon}}
\newunicodechar{δ}{\ensuremath{\delta}}
\newunicodechar{Ε}{\ensuremath{\mathrm{E}}}
\newunicodechar{↔}{\ensuremath{\leftrightarrow}}
\directlua{luaotfload.add_fallback('versomonofallback',
  {'DejaVuSans:mode=harf;', 'STIXTwoMath-Regular:mode=harf;'})}
\setmonofont{DejaVu Sans Mono}[RawFeature={fallback=versomonofallback}]
\RecustomVerbatimCommand{\LeanVerb}{Verb}
  {commandchars=\\\{\},fontsize=\small,breaklines=true,
   breakafter={_./},breakaftersymbolpre={},breakanywhere=true,
   breakanywheresymbolpre={},breaksymbolleft={},breaksymbolright={}}
\emergencystretch=1em
% ---
\begin{document}

\frontmatter

\begin{titlingpage}
\maketitle
\end{titlingpage}

\tableofcontents


\cleardoublepage
A thesis submitted in partial fulfilment of the requirements for the degree of
Doctor of Philosophy.

\mainmatter
\chapter{Abstract}\label{Capability--Monotonicity--in--a--Formal--Skeleton--for--Agentic--AI--Safety----Abstract}

An earlier draft of this project proposed a full dynamic-programming theory
of agentic AI safety — a robust reachability value, a viability kernel, a
critical surface, a recoverability boundary, an adversarial-energy metric,
a memory-hysteresis prediction, and an eight-item empirical program built
on top of them. This document is not that draft. It is a rewrite under a
rule borrowed from its companion thesis, \textbf{Two Schemas of Impossibility}:
a claim appears only if it is proved in an elaborated Lean block or
resolves to a named Lean declaration: it is a thesis about what remains of
a larger proposal once every claim without a Lean declaration behind it is
removed.

What remains is a small formal skeleton — an augmented agent state, a
belief interface, an intervention/composition mechanism, a capability
partial order, a specification automaton, a shield interface, and a
recoverability-monitor interface, each a Lean \LeanVerb|structure|, \LeanVerb|inductive|, or
\LeanVerb|def| checked by the compiler as this document builds — together with one
genuine theorem, \LeanVerb|capability_reachable_subset|: a capability ordering that
only ever adds executable actions can only ever enlarge the reachable
state graph. The theorem is a general fact about monotone reachability
graphs; its relevance to agentic AI safety is a reading argued for in the
text, not a further mathematical result.

Chapter 7 lists every claim that was cut to reach this document — the
value function, the viability kernel, the critical surface, the tail-
severity and adversarial-energy functionals, the recoverability
probability and its hysteresis prediction, a runtime-assurance design
rule, a causal-incentive reading, and the empirical identification
program — and says why each was cut: none of them corresponded to a Lean
declaration. This thesis's claim about itself is only that the skeleton
and the one theorem in Chapters 3 through 6 are genuinely backed, and that
everything listed in Chapter 7 is a research proposal, not a result.




\chapter{How to read this thesis}\label{Capability--Monotonicity--in--a--Formal--Skeleton--for--Agentic--AI--Safety----How--to--read--this--thesis}

This is a short document by design. Chapter 1 states the scope reduction
directly. Chapter 2 places the surviving formal objects, briefly, against
their literatures. Chapters 3 through 6 give the Lean content: the state
and capability structures, the specification automaton, the shield
interface and the capability-monotonicity theorem, and the
recoverability-monitor interface. Chapter 7 is an explicit accounting of
what was removed and why. Chapter 8 concludes.




\chapter{Introduction}\label{Capability--Monotonicity--in--a--Formal--Skeleton--for--Agentic--AI--Safety----Introduction}

This thesis started from a much larger proposal: a full dynamic-programming
theory of agentic AI safety, with a robust reachability value, a viability
kernel, a critical surface, a recoverability boundary, an adversarial-energy
metric, a hysteresis prediction, and an eight-point empirical program. That
proposal is not what follows. It is not what follows because none of it was
backed by a machine-checked declaration, and this thesis's governing rule,
inherited from its companion volume \textbf{Two Schemas of Impossibility}, is that a
claim without Lean backing does not appear in the document. The companion
thesis enforces this with a build-time citation gate: every numbered claim
must either be proved in an elaborated code block or resolve to a named
declaration in one of its four companion libraries, on pain of the build
failing. This thesis adopts the same rule and applies it to itself, and the
result is a much smaller document than the one first proposed.



\section{What survives the rule}\label{Capability--Monotonicity--in--a--Formal--Skeleton--for--Agentic--AI--Safety----Introduction----What--survives--the--rule}

Four objects survive. Chapter 3 gives the augmented agent state, a belief
interface, the intervention/composition mechanism, and the capability
partial order, each as a Lean \LeanVerb|structure|, \LeanVerb|inductive|, or \LeanVerb|def| that the
compiler checks for well-formedness as this document builds. Chapter 4 gives
the specification automaton that a trajectory-level safety property compiles
into. Chapter 5 gives a shield interface and proves, as a genuine
machine-checked theorem, that a capability ordering which only ever adds
executable actions can only ever enlarge what is reachable — not a deep
result, but a real one, \LeanVerb|capability_reachable_subset|, checked by the Lean
kernel rather than asserted in prose. Chapter 6 gives a recoverability
monitor interface, offered explicitly as bookkeeping rather than as backing
for any claim about what recoverability \textbf{is}.



\section{What does not survive}\label{Capability--Monotonicity--in--a--Formal--Skeleton--for--Agentic--AI--Safety----Introduction----What--does--not--survive}

Chapter 7 lists what was cut and why, so that a reader who saw an earlier,
larger version of this project — or who read the informal writeup this
thesis grew out of — can see exactly where the boundary was drawn and can
check that the cut was not selective. The short version: a robust failure
value, a viability kernel, a critical surface, a tail-severity functional, an
adversarial-energy metric, a hysteresis prediction, a runtime-assurance
fallback rule, a reading of causal-incentive analysis, and an eight-item
empirical program were all proposed at one point and all removed, because
none of them was a Lean declaration and this thesis's rule does not make an
exception for an idea being interesting.



\section{Thesis statement, scaled to what remains}\label{Capability--Monotonicity--in--a--Formal--Skeleton--for--Agentic--AI--Safety----Introduction----Thesis--statement-005F-005F-005F--scaled--to--what--remains}

\textbf{A small fragment of agentic-AI-safety reasoning — the state an agent
occupies, the actions available to it under a capability ordering, the
automaton a trajectory-level specification compiles into, and the fact that
enlarging capability can only enlarge what is reachable — can be given a
machine-checked formal encoding in Lean 4, and exactly one nontrivial
consequence of that encoding, capability monotonicity of reachability, is
proved rather than asserted.}

This is a narrower claim than the one the larger proposal made, and it is
narrower on purpose. A reader who wants the larger theory — the dynamic
game, the critical surface, the empirical program — should read it as a
research proposal, stated as such, and should not mistake anything in
Chapters 3 through 6 of this document for evidence toward it.



\section{Structure}\label{Capability--Monotonicity--in--a--Formal--Skeleton--for--Agentic--AI--Safety----Introduction----Structure}

Chapter 2 places the surviving formal objects against their literature.
Chapter 3 gives the state, belief, intervention, and capability structures.
Chapter 4 gives the specification automaton. Chapter 5 gives the shield
interface and the capability-monotonicity theorem. Chapter 6 gives the
recoverability-monitor interface. Chapter 7 lists what was cut, and why.
Chapter 8 concludes.





\chapter{Related Work}\label{Capability--Monotonicity--in--a--Formal--Skeleton--for--Agentic--AI--Safety----Related--Work}

The formal objects that survive into Chapters 3 through 6 are small
instances of larger, established constructions. This chapter places each
one, and is explicit that placing an object against a literature is not the
same as claiming this thesis reproduces that literature's theorems: the
Lean content of this thesis is the four structures and the one theorem
listed in Chapter 1, and nothing said here should be read as extending that
content.



\section{Hybrid state and partial observability}\label{Capability--Monotonicity--in--a--Formal--Skeleton--for--Agentic--AI--Safety----Related--Work----Hybrid--state--and--partial--observability}

Chapter 3's augmented state, with continuous and discrete components living
side by side, is an instance of the standard object in stochastic hybrid
systems theory, where reachability results are already established for
systems combining both. Chapter 3's belief interface is the standard
sufficient-statistic construction from partially observable planning: the
posterior over the hidden state given the observation history is what a
planner needs to act optimally under partial observability, and belief-space
planning is the classical way to reduce a partially observed problem to a
fully observed one over that posterior. This thesis's \LeanVerb|Belief| structure
captures only the interface — a point estimate and an update rule — not the
measure-theoretic content, deliberately: nothing downstream in this thesis
needs more than the interface.



\section{Specification automata}\label{Capability--Monotonicity--in--a--Formal--Skeleton--for--Agentic--AI--Safety----Related--Work----Specification--automata}

Compiling a trajectory-level temporal property into a monitor automaton,
which Chapter 4's \LeanVerb|SpecAutomaton| does, is the standard construction from
runtime verification and from shielding in safe reinforcement learning: a
temporal-logic specification is compiled into an automaton whose state
tracks progress toward, or violation of, the property, so that a
trajectory-level question becomes an ordinary question about the automaton's
current state. Shielding then overrides a learner's action only when doing
otherwise would violate the specification, which is the idea Chapter 5's
\LeanVerb|Shield| structure encodes, again as an interface rather than as a proof
that any concrete shield built on it satisfies the specification.



\section{Capability and reachability}\label{Capability--Monotonicity--in--a--Formal--Skeleton--for--Agentic--AI--Safety----Related--Work----Capability--and--reachability}

The claim proved in Chapter 5 — that a capability ordering which only adds
executable actions can only enlarge the reachable set — is an instance of an
entirely general fact about reachability in monotone systems, unrelated to
agents specifically: it holds for any state-transition graph where a
"stronger" edge set is a superset of a "weaker" one. Its relevance to
agentic AI safety is the reading, not the mathematics: shell access,
unrestricted networking, credentials, write access, and delegation are each
naturally modeled as enlarging the executable-action set at every state that
has them, and the theorem then says the resulting reachable set can only
grow. This reading is argued for in Chapter 5; the theorem itself needs
nothing about agents to be true.



\section{What this chapter does not place}\label{Capability--Monotonicity--in--a--Formal--Skeleton--for--Agentic--AI--Safety----Related--Work----What--this--chapter--does--not--place}

Belief-space robustness under model uncertainty, temporal reachability
values, runtime-assurance switching, information-flow reference monitors,
and causal-incentive analysis all have substantial literatures that an
earlier draft of this thesis related its (since-removed) claims to. Chapter
7 records that those claims were removed and does not restore this
chapter's discussion of their literatures, because a related-work section
for a claim that no longer appears in the thesis would misrepresent what the
thesis contains.





\chapter{The Formal Model}\label{Capability--Monotonicity--in--a--Formal--Skeleton--for--Agentic--AI--Safety----The--Formal--Model}

This chapter fixes the objects that are actually declared in Lean and
checked by the compiler as this document builds. Each structure below is
followed immediately by the code that defines it; nothing is described here
that is not in that code.



\section{The augmented state}\label{Capability--Monotonicity--in--a--Formal--Skeleton--for--Agentic--AI--Safety----The--Formal--Model----The--augmented--state}

The state of a closed-loop agentic system, at a point in time, is modeled as
a tuple of a latent computational component, a discrete operating mode, a
memory/context/provenance component, a capability component, an
environment component, and the state of a specification monitor (Chapter 4).

\begin{LeanVerbatim}
\textbf{structure} CTVGState (\textit{X} \textit{Q} \textit{M} \textit{C} \textit{E} \textit{RPhi} : Type) \textbf{where}
  latent      : \textit{X}
  mode        : \textit{Q}
  memory      : \textit{M}
  capability  : \textit{C}
  environment : \textit{E}
  monitor     : \textit{RPhi}

\end{LeanVerbatim}


This is a product type and nothing more: the compiler checks that it forms
a valid structure, which is a weaker guarantee than any theorem about it,
and no theorem about \LeanVerb|CTVGState| is proved anywhere in this thesis.



\section{Belief as an interface}\label{Capability--Monotonicity--in--a--Formal--Skeleton--for--Agentic--AI--Safety----The--Formal--Model----Belief--as--an--interface}

A belief is modeled here as an opaque summary type together with the two
operations a monitor needs: reading off a point estimate, and updating on
an observation.

\begin{LeanVerbatim}
\textbf{structure} Belief (\textit{S} \textit{Y} : Type) \textbf{where}
  estimate : \textit{S}
  update   : \textit{Y} → \textit{S}

\end{LeanVerbatim}


This is deliberately not a probability measure. A concrete instance
(particle filter, Gaussian belief, finite-support posterior) would refine
\LeanVerb|estimate| and \LeanVerb|update|; nothing in this thesis provides such an instance,
and no claim here depends on which refinement a future instance would
choose.



\section{Intervention and composition}\label{Capability--Monotonicity--in--a--Formal--Skeleton--for--Agentic--AI--Safety----The--Formal--Model----Intervention--and--composition}

An agent proposes an action; a defender's intervention either lets it
through or replaces it.

\begin{LeanVerbatim}
\textbf{inductive} Intervention (\textit{A} : Type) \textbf{where}
  \symbol{124} allow
  \symbol{124} modify (\textit{a'} : \textit{A})
  \symbol{124} sandbox (\textit{a'} : \textit{A})
  \symbol{124} reducePrivilege (\textit{a'} : \textit{A})
  \symbol{124} askHuman
  \symbol{124} deny
  \symbol{124} fallback

\textbf{def} compose \symbol{123}\textit{A} : Type\symbol{125} (\textit{proposed} : \textit{A}) : Intervention \textit{A} → \textit{A}
  \symbol{124} .allow              => \textit{proposed}
  \symbol{124} .modify \textit{a'}          => \textit{a'}
  \symbol{124} .sandbox \textit{a'}         => \textit{a'}
  \symbol{124} .reducePrivilege \textit{a'} => \textit{a'}
  \symbol{124} .askHuman           => \textit{proposed}
  \symbol{124} .deny               => \textit{proposed}
  \symbol{124} .fallback           => \textit{proposed}

\end{LeanVerbatim}


\LeanVerb|compose| is total and type-checks; that is the entire content of this
section. The three cases that currently fall through to \LeanVerb|proposed|
(\LeanVerb|askHuman|, \LeanVerb|deny|, \LeanVerb|fallback|) do so because no distinguished "safe
default" action has been fixed at this level of generality — fixing one
would require committing to a concrete \LeanVerb|Action| type, which this chapter
does not do.



\section{Capability ordering}\label{Capability--Monotonicity--in--a--Formal--Skeleton--for--Agentic--AI--Safety----The--Formal--Model----Capability--ordering}

Capability sets are given a partial order: \LeanVerb|c1 ≤ c2| means every action
\LeanVerb|c1| allows, \LeanVerb|c2| also allows.

\begin{LeanVerbatim}
\textbf{structure} CapabilitySet (\textit{Action} : Type) \textbf{where}
  executable : \textit{Action} → Bool

\textbf{def} CapabilitySet.le \symbol{123}\textit{Action} : Type\symbol{125} (\textit{c1} \textit{c2} : CapabilitySet \textit{Action}) : Prop :=
  ∀ \textit{a}, \textit{c1}.executable \textit{a} = true → \textit{c2}.executable \textit{a} = true

\textbf{instance} \symbol{123}\textit{Action} : Type\symbol{125} : LE (CapabilitySet \textit{Action}) := ⟨CapabilitySet.le⟩

\end{LeanVerbatim}


This ordering is what Chapter 5's theorem is stated against. It is a
\LeanVerb|def| and an \LeanVerb|instance|, checked for type-correctness by the compiler; the
theorem that gives this ordering its content is proved separately, in
Chapter 5, and is not restated here.





\chapter{The Specification Automaton}\label{Capability--Monotonicity--in--a--Formal--Skeleton--for--Agentic--AI--Safety----The--Specification--Automaton}

Not every hazard is a property of a single state. Reading an
attacker-controlled document, extracting a secret, and transmitting it to
an attacker-controlled endpoint violates no single-state predicate; the
violation is a property of the trajectory. This chapter gives the one Lean
object this thesis uses to reduce that kind of trajectory-level question to
an ordinary state-level one.



\section{The automaton}\label{Capability--Monotonicity--in--a--Formal--Skeleton--for--Agentic--AI--Safety----The--Specification--Automaton----The--automaton}

A safety specification is compiled into a deterministic monitor with a
start state, a step function reading the current state, the action taken,
and the resulting state, and a predicate marking which of its states are
rejecting.

\begin{LeanVerbatim}
\textbf{structure} SpecAutomaton (\textit{R} \textit{S} \textit{A} : Type) \textbf{where}
  init    : \textit{R}
  step    : \textit{R} → \textit{S} → \textit{A} → \textit{S} → \textit{R}
  reject  : \textit{R} → Bool

/-- Running the automaton over a whole trajectory, so that a trajectory
violates the specification exactly when the automaton ends in a rejecting
state. -/
\textbf{def} SpecAutomaton.run \symbol{123}\textit{R} \textit{S} \textit{A} : Type\symbol{125} (\textit{spec} : SpecAutomaton \textit{R} \textit{S} \textit{A})
    (\textit{states} : List \textit{S}) (\textit{actions} : List \textit{A}) : \textit{R} :=
  \textbf{match} \textit{states}, \textit{actions} \textbf{with}
  \symbol{124} \textit{s} :: \textit{s'} :: \textit{states'}, \textit{a} :: \textit{actions'} =>
      \textit{spec}.step (SpecAutomaton.run \textit{spec} (\textit{s} :: \textit{states'}) \textit{actions'}) \textit{s} \textit{a} \textit{s'}
  \symbol{124} _, _ => \textit{spec}.init

\end{LeanVerbatim}


\LeanVerb|SpecAutomaton.run| type-checks and its termination is accepted by the
compiler's structural-recursion check on the argument lists; nothing beyond
that is claimed. In particular, this chapter does not prove that any
concrete instantiation of \LeanVerb|step| correctly compiles any particular temporal
formula, does not prove \LeanVerb|SpecAutomaton.run| correct against any semantics
for trajectories, and does not prove that \LeanVerb|reject| classifies trajectories
the way an informal safety specification would want. What is proved is only
that the automaton, as a piece of Lean code, is well-typed and terminates —
which is exactly the reduction Chapter 3's \LeanVerb|CTVGState.monitor| field
presupposes is available: the monitor component of the state is intended to
be populated by \LeanVerb|SpecAutomaton.run|'s current partial result, so that a
trajectory-level property becomes the ordinary state-level question of
whether that field currently satisfies \LeanVerb|reject|.





\chapter{The Shield, and Capability Monotonicity}\label{Capability--Monotonicity--in--a--Formal--Skeleton--for--Agentic--AI--Safety----The--Shield-005F-005F-005F--and--Capability--Monotonicity}

This chapter contains the one theorem in this thesis that is a genuine
machine-checked proof rather than a definition, an interface, or a
conjecture.



\section{The shield interface}\label{Capability--Monotonicity--in--a--Formal--Skeleton--for--Agentic--AI--Safety----The--Shield-005F-005F-005F--and--Capability--Monotonicity----The--shield--interface}

A shield checks an agent's proposed action against a hard admissibility
predicate and a separate, second-stage predicate, and substitutes the best
available admissible alternative when the proposal fails either check.

\begin{LeanVerbatim}
/-- A shield interface: `hardInvariant` and `valueBound` are left abstract,
so that this structure fixes only the control flow — check, then fall back
to the best admissible alternative — not any concrete safety criterion. -/
\textbf{structure} Shield (\textit{Action} : Type) \textbf{where}
  hardInvariant : \textit{Action} → Bool
  valueBound    : \textit{Action} → Bool
  utility       : \textit{Action} → Float
  candidates    : List \textit{Action}

\textbf{def} Shield.admissible \symbol{123}\textit{Action} : Type\symbol{125} (\textit{shield} : Shield \textit{Action}) : List \textit{Action} :=
  \textit{shield}.candidates.filter (\textbf{fun} \textit{a} => \textit{shield}.hardInvariant \textit{a} && \textit{shield}.valueBound \textit{a})

\textbf{def} Shield.decide \symbol{123}\textit{Action} : Type\symbol{125} (\textit{shield} : Shield \textit{Action}) (\textit{proposed} : \textit{Action})
    (\textit{isProposedAdmissible} : Bool) : \textit{Action} :=
  \textbf{if} \textit{isProposedAdmissible} \textbf{then} \textit{proposed}
  \textbf{else}
    \textbf{let} \textit{admissible} := \textit{shield}.admissible
    \textbf{match} \textit{admissible} \textbf{with}
    \symbol{124} []      => \textit{proposed}  -- no fallback is fixed at this level of generality
    \symbol{124} \textit{a} :: \textit{as} =>
      \textit{as}.foldl (\textbf{fun} \textit{best} \textit{a} =>
        \textbf{if} \textit{shield}.utility \textit{a} > \textit{shield}.utility \textit{best} \textbf{then} \textit{a} \textbf{else} \textit{best}) \textit{a}

\end{LeanVerbatim}


\LeanVerb|Shield.decide|'s empty-admissible-set branch returns the unmodified
proposal, which is very likely the wrong thing to do in a deployed system.
No theorem in this thesis says what the right thing is, and no default is
fixed here beyond what type-checks.



\section{Capability monotonicity}\label{Capability--Monotonicity--in--a--Formal--Skeleton--for--Agentic--AI--Safety----The--Shield-005F-005F-005F--and--Capability--Monotonicity----Capability--monotonicity}

Reachability in a finite number of steps, under a given capability set, is
defined as an inductive predicate: a state reaches itself in zero steps, and
reaches a further state by one more executable-action step.

\begin{LeanVerbatim}
\textbf{inductive} ReachableIn \symbol{123}\textit{State} \textit{Action} : Type\symbol{125} (\textit{step} : \textit{State} → \textit{Action} → List \textit{State})
    (\textit{c} : CapabilitySet \textit{Action}) : \textit{State} → \textit{State} → Prop \textbf{where}
  \symbol{124} refl (\textit{s} : \textit{State}) : ReachableIn \textit{step} \textit{c} \textit{s} \textit{s}
  \symbol{124} tail \symbol{123}\textit{s0} \textit{s'} \textit{s} : \textit{State}\symbol{125} (\textit{a} : \textit{Action}) :
      ReachableIn \textit{step} \textit{c} \textit{s0} \textit{s'} →
      \textit{c}.executable \textit{a} = true →
      \textit{s} ∈ \textit{step} \textit{s'} \textit{a} →
      ReachableIn \textit{step} \textit{c} \textit{s0} \textit{s}

\textbf{theorem} capability_reachable_subset \symbol{123}\textit{State} \textit{Action} : Type\symbol{125}
    (\textit{step} : \textit{State} → \textit{Action} → List \textit{State})
    (\textit{c1} \textit{c2} : CapabilitySet \textit{Action}) (\textit{h} : \textit{c1}.le \textit{c2})
    (\textit{s0} \textit{s} : \textit{State}) :
    ReachableIn \textit{step} \textit{c1} \textit{s0} \textit{s} → ReachableIn \textit{step} \textit{c2} \textit{s0} \textit{s} := \textbf{by}
  \textbf{intro} \textit{hr}
  \textbf{induction} \textit{hr} \textbf{with}
  \symbol{124} refl => \textbf{exact} ReachableIn.refl \textit{s0}
  \symbol{124} tail \textit{a} \textit{hprev} \textit{hexec} \textit{hstep} \textit{ih} =>
      \textbf{exact} ReachableIn.tail \textit{a} \textit{ih} (\textit{h} \textit{a} \textit{hexec}) \textit{hstep}

\end{LeanVerbatim}


\LeanVerb|capability_reachable_subset| is checked by the Lean kernel: if \LeanVerb|c1 ≤ c2|
in the sense of Chapter 3's capability ordering, every state reachable
under \LeanVerb|c1| is reachable under \LeanVerb|c2|. Its content is exactly this statement
and no more. It does not say anything about probability, about a defender's
ability to control what is reached, about safety, or about usefulness; it
is a fact about monotone reachability graphs, stated here because the
capability ordering it needs was already in scope, and its relevance to
agentic AI safety is entirely in the reading — shell access, unrestricted
networking, credentials, write access, and delegation are each naturally
modeled as enlarging the executable-action set at every state that has
them — a reading this thesis asserts but does not, and could not, prove.





\chapter{A Recoverability-Monitor Interface}\label{Capability--Monotonicity--in--a--Formal--Skeleton--for--Agentic--AI--Safety----A--Recoverability--Monitor--Interface}

An earlier draft of this thesis defined recoverability as a probability —
the best-case-over-defender, worst-case-over-adversary chance of returning
to a confidently safe core region before failing — and asserted, without
proof, that crossing a resulting boundary was the event most worth
monitoring in a deployed agent. That probability was never computed, no
existence result was given for it, and no Lean object backed it, so under
this thesis's rule it was removed (Chapter 7). What remains is the
interface a monitor implementing some future, actual recoverability
estimate would need.

\begin{LeanVerbatim}
/-- An interface for a recoverability monitor: given a current estimate and
an observation, produce the next estimate, and flag the event that the
estimate has dropped from at or above a threshold to below it. Nothing here
computes, bounds, or justifies `estimate`; the structure only fixes what a
caller would pass in and get back. -/
\textbf{structure} RecoverabilityMonitor (\textit{Y} : Type) \textbf{where}
  estimate : Float
  update   : \textit{Y} → Float

\textbf{def} crossedPointOfNoReturn (\textit{eta} \textit{prevEstimate} \textit{newEstimate} : Float) : Bool :=
  (1.0 - \textit{eta} ≤ \textit{prevEstimate}) && (\textit{newEstimate} < 1.0 - \textit{eta})

\end{LeanVerbatim}


\LeanVerb|RecoverabilityMonitor| and \LeanVerb|crossedPointOfNoReturn| type-check; that is the
whole of this chapter's content. \LeanVerb|crossedPointOfNoReturn| is a pure
comparison of two floating-point numbers against a threshold — it says
nothing about what \LeanVerb|estimate| should be computed from, and no instance of
\LeanVerb|RecoverabilityMonitor| is constructed anywhere in this thesis. The chapter
is included, rather than cut along with the rest of the recoverability
material, because it is genuinely Lean: two declarations that type-check
against Chapter 3's structures and nothing more.





\chapter{What Was Cut, and Why}\label{Capability--Monotonicity--in--a--Formal--Skeleton--for--Agentic--AI--Safety----What--Was--Cut-005F-005F-005F--and--Why}

This thesis was rewritten under a rule adopted from its companion volume:
a claim survives only if it is proved in an elaborated Lean block or
resolves to a named declaration. Applying that rule to the original
proposal removed most of it. This chapter lists what was removed, so the
cut is auditable rather than quiet.



\section{The robust failure value and the viability kernel}\label{Capability--Monotonicity--in--a--Formal--Skeleton--for--Agentic--AI--Safety----What--Was--Cut-005F-005F-005F--and--Why----The--robust--failure--value--and--the--viability--kernel}

Proposed: a value function \LeanVerb|V_T(B)|, computed as the smallest,
over defender policies, of the largest, over adversary policies and
plausible transition models, probability of a trajectory reaching a
failure set before a completion set; a "useful viability kernel" of belief
states from which some defender policy keeps this value and a companion
tail-severity functional below fixed tolerances while still finishing the
task; and a "critical surface" at the kernel's boundary.

Cut, in full. No transition model, no defender-policy type, no
adversary-policy type, and no value function were ever defined in Lean.
The recursion proposed to compute \LeanVerb|V_T| was written down by analogy with
classical Bellman recursions and its existence and convergence were never
established, even informally, let alone proved.



\section{Tail severity and adversarial energy}\label{Capability--Monotonicity--in--a--Formal--Skeleton--for--Agentic--AI--Safety----What--Was--Cut-005F-005F-005F--and--Why----Tail--severity--and--adversarial--energy}

Proposed: a CVaR-based tail-severity functional, and a "critical energy"
measuring the minimum adversarial effort needed to force a belief state out
of the viability kernel.

Cut, in full, for the same reason: both were defined only in prose, against
objects (the viability kernel, an attacker effort function) that themselves
were never formalized.



\section{Recoverability as a probability, and the hysteresis prediction}\label{Capability--Monotonicity--in--a--Formal--Skeleton--for--Agentic--AI--Safety----What--Was--Cut-005F-005F-005F--and--Why----Recoverability--as--a--probability-005F-005F-005F--and--the--hysteresis--prediction}

Proposed: a recoverability functional \LeanVerb|R_T(B)|, a "core" safe region, a
resulting point-of-no-return boundary, and a prediction that memory-bearing
agents would show measurable hysteresis in failure probability under a
forward/backward sweep of a hazard-intensity parameter, while stateless
agents would not.

Cut. What is kept, in Chapter 6, is only the bookkeeping interface a future
monitor implementing some such quantity would expose — an estimate and an
update function — with no claim about what the estimate should be, how it
would be computed, or whether the hysteresis prediction is true. The
prediction itself, being an empirical claim about systems that were never
built or measured as part of this thesis, does not belong in a document
gated on Lean backing, and is removed rather than retained as a labeled
conjecture.



\section{The runtime-assurance fallback rule}\label{Capability--Monotonicity--in--a--Formal--Skeleton--for--Agentic--AI--Safety----What--Was--Cut-005F-005F-005F--and--Why----The--runtime--assurance--fallback--rule}

Proposed: a design rule that, when a shield's admissible-action set is
empty, control should switch to a distinct, independently-certified
fallback controller, by analogy with the Simplex architecture from
safety-critical control.

Cut as a rule. What is kept, in Chapter 5, is the observation that
\LeanVerb|Shield.decide|'s empty-admissible-set branch currently returns the
unmodified proposal — a fact about the Lean code, stated as a fact about
the code, not as a design recommendation about what a real system should
do instead.



\section{Causal-incentive analysis and the Defense Trilemma reading}\label{Capability--Monotonicity--in--a--Formal--Skeleton--for--Agentic--AI--Safety----What--Was--Cut-005F-005F-005F--and--Why----Causal--incentive--analysis--and--the--Defense--Trilemma--reading}

Proposed: a reading of structural causal influence diagrams as
complementary to reachability analysis, identifying "pressure" where
reachability identifies "consequence"; and a reading of the companion
thesis's Defense Trilemma result as motivating this thesis's (now-removed)
shield architecture.

Cut, in full. Neither reading was a claim this thesis could back with a
declaration — the first cites no diagram formalized anywhere in this
project, and the second is commentary on a different thesis's theorem, not
a claim proved here.



\section{The eight-item empirical identification program}\label{Capability--Monotonicity--in--a--Formal--Skeleton--for--Agentic--AI--Safety----What--Was--Cut-005F-005F-005F--and--Why----The--eight--item--empirical--identification--program}

Proposed: a procedure for discovering an adequate state representation by
an empirical Markov-sufficiency test, a procedure for fitting a transition
model and an ambiguity set from data, a procedure for estimating \LeanVerb|V_T| and
testing for a critical surface, and eight falsifiable predictions collected
from the rest of the (now-removed) theory.

Cut, in full. No data was collected, no model was fit, and every prediction
depended on a quantity — \LeanVerb|V_T|, \LeanVerb|E_crit|, \LeanVerb|R_T|, the hysteresis loop area —
that no longer appears in this thesis.



\section{What this leaves}\label{Capability--Monotonicity--in--a--Formal--Skeleton--for--Agentic--AI--Safety----What--Was--Cut-005F-005F-005F--and--Why----What--this--leaves}

What survives is Chapters 3 through 6: four Lean structures, one inductive
predicate, and one theorem. That is the whole of this thesis's formal
content, and Chapter 8 says so without padding it back out.





\chapter{Conclusion}\label{Capability--Monotonicity--in--a--Formal--Skeleton--for--Agentic--AI--Safety----Conclusion}

This thesis, in its current form, formalizes a small piece of vocabulary
for reasoning about agentic AI safety — an augmented state, a belief
interface, an intervention/composition mechanism, a capability partial
order, a specification automaton, a shield interface, and a
recoverability-monitor interface — and proves one theorem against that
vocabulary: \LeanVerb|capability_reachable_subset|, that a capability ordering which
only ever adds executable actions can only ever enlarge what is reachable.

That is a modest result relative to the theory originally proposed. It is
stated as such rather than dressed up, because the alternative — keeping
the larger theory's definitions and conjectures in the document, labeled as
conjectures — was tried in an earlier draft and rejected: a claim that
cannot be checked by the Lean kernel does not belong in a thesis gated on
Lean backing, no matter how carefully it is hedged in prose. Chapter 7
records exactly what that rejection removed, so that the removal is
inspectable rather than a matter of trust.

What this thesis takes itself to have shown is narrower than what its
companion volume shows for hallucination and prompt-injection defense: a
type-correct formal skeleton, and one theorem within it. The larger claims
about robust reachability values, viability kernels, critical surfaces,
recoverability, adversarial energy, and hysteresis in memory-bearing agents
remain exactly what they were before this rewrite — informal proposals,
possibly worth formalizing — and this document no longer asserts them.





\end{document}
